\begin{textsample}{POS Dim 3 – pi – Score 6.00 – pi/pi\_ef\_2.txt}  \label{ex:f3_pos_268}
OBJETIVOS DO COMPONENTE LÍNGUA PORTUGUESA A PARTIR DA BASE NACIONAL COMUM CURRICULAR Neste currículo, valorizamos a natureza discursiva da língua e compreendemos que ela produz efeitos de sentido ( PÊCHEUX, 2010 ) sobre um objeto, pois a elaboração de sentido depende não apenas da materialidade da língua, mas também da exterioridade que a permeia.
Nessa exterioridade, encontram -se os sujeitos que agem na sociedade, a história, condições de produção e os próprios saberes sociais ( ORLANDI, 2012 ).
Nesse ínterim, é imprescindível que haja formação tanto no nível de alfabetização como no que diz respeito ao letramento.
Alfabetização e Letramento A alfabetização ocorre com a codificação e a decodificação do sistema de escrita alfabética.
Assim, ela envolve o desenvolvimento de uma consciência fonológica ( dos fonemas do português do Brasil e de sua organização em segmentos sonoros maiores como sílabas e palavras ) quer seja em formato cursivo, quer seja em imprensa, bem como em outras apresentações possíveis.
Logo, esse processo permeia a compreensão dos grafemas e fonemas de palavras em textos escritos.
De acordo com a BNCC, é preciso que ela seja implementada nos dois primeiros anos, e deve ser o objetivo principal da ação pedagógica deste componente.
O Letramento, por sua vez, é o processo pelo qual se faz uso competente da leitura e da escrita em situações reais, em que se deve exceder as ações de codificar e decodificar o sistema.
Tal processo requer metodologias que tenham significados para o aluno e levem -no a refletir sobre o emprego da língua nas práticas sociais.
Ele deve se fazer presente desde as esferas primárias, mais simples, em situações cotidianas, como leitura de panfletos ou placas, até as secundárias, mais complexas, em situações mais específicas, como leitura de textos jurídicos ou científicos.
No que diz respeito às maneiras de alfabetizar tomando como base o letramento, propõe -se o direcionamento de práticas que partam da leitura dos textos e cheguem às práticas educativas que explorem a linguagem e o universo \textbf{infantil} \textbf{lúdico} e das \textbf{brincadeiras}.
Isso deve resultar na ampliação da visão de mundo e das possibilidades de se utilizar a língua nas diversas situações cotidianas.
É necessário, pois, proporcionar ao estudante piauiense oportunidades para se expressar, utilizando -se das práticas de linguagem em situações reais.
Nesse sentido, a leitura acontece em um mútuo e tácito acordo entre materialidade linguística e exterioridade, ambas importantes.
Esses pensamentos nos levaram a refletir sobre como deveríamos considerar elementos fundamentais e necessários para o ensino de Língua Portuguesa.
Concluímos que o texto é um material que tem em si diferentes discursos e que ativa, para sua produção, leitura ou \textbf{escuta}, conhecimentos diferentes.
Isso explica por que o tomamos como basilar para o ensino.
É uma mediação entre sujeitos ( ORLANDI, 2012b ).
Em vista disso, tencionamos que os objetivos postos neste documento promovam uma efetiva participação dos indivíduos como cidadãos, sujeitos de direitos e deveres, por meio de textos em gêneros diversos.
O texto aqui também se assume como um objeto de análise capaz de promover inter-relações com outros, verbais e/ou não verbais, e outras artes como música, filmes, esculturas, \textbf{pinturas}.
Acreditamos que existam inúmeras formas de textos, dependendo da sua criação.
Não julgamos como produção textual apenas aquela que utiliza palavras como matéria-prima, mas também as que utilizam imagens, \textbf{sons}, \textbf{gestos}, o que configura textos multimodais.
Essa forma de leitura, ao promover a aproximação com tantos meios, lembra -nos o Currículo do Piauí, que, por meio de suas proposições, permite aos professores trabalhar de forma interdisciplinar em todos os anos do ensino fundamental, com início no primeiro ano, evitando a fragmentação de conteúdos e propiciando o diálogo entre os componentes curriculares.
Em articulação com as competências gerais da Educação Básica e com as competências específicas da área de Linguagens, o componente curricular de Língua Portuguesa deve garantir aos estudantes piauienses o desenvolvimento de competências específicas, sejam elas cognitivas ou socioemocionais.
Vale ainda destacar que as competências perpassam todos os componentes curriculares do Ensino Fundamental e são essenciais para a ampliação das possibilidades de participação dos estudantes em práticas de diferentes campos de atividades humanas e de pleno exercício da cidadania.
Temos as seguintes competências específicas que nos ajudam na elaboração deste currículo, na perspectiva de educação integral.

% matched lemmas: brincadeira, escuta, gesto, infantil, lúdico, pintura, som
\end{textsample}
